\documentclass{article}
\usepackage[T1]{fontenc}
\usepackage{lmodern}
\usepackage{graphicx}
\usepackage{amsmath,amssymb}
\usepackage[a4paper,margin=2cm]{geometry}
\usepackage[absolute,overlay]{textpos}
\setlength{\TPHorizModule}{1mm}
\setlength{\TPVertModule}{1mm}
\usepackage[indonesian]{babel}
\usepackage{hyperref}
\setlength{\emergencystretch}{2em}
% unit: Halaman 34

\begin{document}

% === Halaman 34 (llm) ===

2. OAEP (\textit{Optimal Asymmetric Encryption Padding})
\begin{itemize}
    \item OAEP adalah metode \textit{padding} yang ditambahkan ke pesan sebelum dienkripsi dengan RSA.
    \item \textit{Padding} diperlukan karena:
    \begin{itemize}
        \item RSA tidak boleh mengenkripsi pesan langsung
        \item Tanpa \textit{padding}, RSA rentan terhadap berbagai serangan kriptografi
    \end{itemize}
    \item OAEP membuat pesan menjadi acak (\textit{randomized}) sebelum dienkripsi RSA, sehingga aman terhadap \textit{chosen-ciphertext attack}
    \begin{itemize}
        \item Skema ini diperkenalkan oleh Mihir Bellare dan Phillip Rogaway pada tahun 1994.
    \end{itemize}
    \item Pesan M sebelum dienkripsi RSA diproses sebagai berikut:
    \[ M \rightarrow OAEP\ padding \rightarrow encoded\ message \rightarrow RSA\ encryption \]
    \item Sehingga \textit{ciphertext} menjadi tidak deterministik (setiap enkripsi berbeda walaupun pesannya sama).
\end{itemize}

\end{document}
